

\chapter{Basisuntersuchungen}

% \section{Untersuchungen können uns weiterhelfen}





% \section{Blutdruck- und Herzfrequenzmessung}

\label{topic:rr-messung}

% Grundlagen
% 
% Fallstricke
% 
% Die Blutdruckmanschette soll nach Möglichkeit nicht auf der Seite
% 
% \begin{itemize}
%     \item eines Shuntarmes bei Dialysepatienten
%     \begin{itemize}
%         \item Ein Dialyse-Shunt ist ein künstlicher Kurzschluss zwischen einer
%         Armarterie und einer Armvene. Dadurch wird die Vene erweitert und sie
%         pulsiert, mann kann viel mehr Durchfluss für die Dialyse erreichen.
%         Die Nahtstelle zwischen Vene und Arterie ist jedoch sehr sensibel, wenn
%         gestaut wird kann sie aufplatzen 
%         \item ${\rightarrow}$ Gefahr einer massiven Shuntblutung
%     \end{itemize}
%     \item eines Armödems
%     \item einer Brust-OP mit Lymphknoten-Entfernung (\"Odemgefahr)
%     \item einer Halbseitenschwäche oder -lähmung
% \end{itemize}
% 
% angelegt werden. 
% 
% \gCave{Shunt-Arme sind tabu für die Blutdruckmessung!}
% 
% Bitten Sie im Praktikum am KTW Dialysepatienten höflich Ihnen ihren
% Shunt-Arm zu zeigen!
% 
% Maßnahmen
% 
% \begin{enumerate}
%     \item Evt. Funktion des Blutdruckmessgerätes und des Stethoskops prüfen
%     \begin{itemize}
%         \item Kontrolle der Dichte des Ventils, der Verbindungsschläuche und
%         der Manschette (Ventil schließen, Manschette etwas aufpumpen und
%         Zusammenpressen, Manschette ganz entlüften)
%         \item Kontrolle der Gängigkeit des Ventils (es soll sich leicht auf
%         und zuschrauben lassen)
%         \item Kontrolle des Manometers auf korrekte Eichung: Die Nadel muss sich
%         im Ruhezustand genau bei 0\,mmHg befinden
%         \item Ohrstöpsel des Stethoskops zeigen schräg nach vorne
%         \item Trichterfunktion durch Berühren der Membran prüfen
%     \end{itemize}
%     \item Patientenlagerung
%     \begin{itemize}
%         \item Arm freimachen und entspannt auf einen Tisch legen in leichter
%         Flexion und Supination
%         \item Oberarm auf Herzhöhe
%     \end{itemize}
%     \item Auswahl der geeigneten Manschettengröße
%     \begin{itemize}
%         \item Manschette laut Herstellerangaben wählen z.B Standardmanschette:
%         Oberarmumfang max. 31 cm
%         \item Messwertverfälschung vermeiden! Der aufblasbare Teil der
%         Manschette soll 80\% des Oberarmumfangs umfassen, die Breite der
%         Manschette mindestens 1/3 des Oberarmumfangs betragen
%         \item Zu kleine Manschetten führen zum Messen zu hoher Werte, zu weite
%         Manschetten können nicht am Arm fixiert werden
%     \end{itemize}
%     \item Handhabung des Stethoskops
%     \item Bestimmen der Anlegestelle für Manschette und der
%     Auskultationsstelle durch Palpation
%     \begin{itemize}
%         \item Arteria brachialis medial in der Ellenbeuge aufsuchen
%     \end{itemize}
%     
%     \item Anlegen der Manschette
%     \begin{itemize}
%         \item Manschette muss auf unbekleideter Haut aufliegen
%         \item mäßig straff
%         \item Abstand zur Ellenbeuge 2,5 cm
%         \item aufblasbaren Teil über A. brachialis zentrieren
%         \item Schläuche nicht im Weg,
%         \item CAVE: Halbseitenschwäche/-lähmung, Shuntarm bei
%         Dialysepatienten, Armödem
%     \end{itemize}
%     \item Puls fühlen beim Aufpumpen
%     \begin{itemize}
%         \item Schnelles Aufpumpen
%         \item Tasten des Radialispulses mit Finger 2-4
%         \item Aufpumpen bis 30 mmHg über jenen Druck, ab dem kein Radialispuls
%         mehr getastet wird
%     \end{itemize}
%     \item Messvorgang
%     \begin{itemize}
%         \item Kopf des Stethoskops auf a. brachialis legen
%         \item Druck langsam ablassen 2 mmHg/sec
%         \item Systolischer Wert: Beginn der turbulenten Korotkoffschen
%         Strömungsgeräusche
%         \item diastolischer Wert: Verschwinden der Korotkoffschen Geräusche
%         \item Wenn Geräusche überhaupt nicht verschwinden (z.B. bei Kindern,
%         Schwangeren etc.), so gilt als diastolischer Wert der Messpunkt, an dem
%         die Geräusche deutlich leiser werden
%         \item Wenn Geräusche verschwunden sind, noch weitere 10mmHg langsam
%         ablassen -- wenn dann keine Geräusche mehr hörbar sind, kann die
%         Manschette zügig entleert werden
%         \item CAVE: auskultatorische Lücke
%     \end{itemize}
%     \item Protokollierung der Blutdruckwerte
% \end{enumerate}












\section{Sehen -- Hören -- Fühlen }
\gLayout{0}{Unsere eigenen Sinne}{
Die Seh-, Hör- und Tastsinne sind bei der körperlichen
Untersuchung unverzichtbar. Zuerst wird der Patient und sein
Körper betrachtet. Wir versuchen uns bereits bekannte Symptome wahrzunehmen. Hierbei können uns viele
augenscheinliche Sachen auffallen: zum Beispiel
\gE{Fettleibigkeit}, \gE{Zyanose}, \gE{Blässe},
\gE{Schweiß}, \gE{Fremdkörper}, \gE{Verletzungen}, \gE{Fehlstellungen},
\ldots

Der nächste Schritt betrifft das Hören: Hier sind zum
Beispiel die Atemgeräusche einzuordnen, die man eventuell schon ohne
Hilfsmittel vernehmen kann (brodelndes Atemgeräusch beim
Lungenödem, pfeifendes Atemgeräusch beim Asthmaanfall,
röcheln bei Atemwegsverlegung, \ldots)

Wenn wir gesehen und gehört haben können wir versuchen zu
fühlen: Wir können den (Radialis-) Puls fühlen, können
(stehende) Hautfalten beurteilen oder das Abdomen abtasten.
}{
\begin{itemize}
  \item Zyanose
  \item Blässe
  \item Schweiß
  \item Körperhaltung
  \item Fremdkörper
  \item Verletzungen, Fehlstellungen
  \item Atemgeräusche
  \item Puls
  \item Stehende Hautfalten
  \item Durchblutung des Nagelbetts
  \item \ldots
\end{itemize}
}{}{}{}




\section{Vitalwerte}
Das erheben von Vitalwerten ist eine Routinetätigkeit und
sollte bei jedem Notfallpatienten erfolgen. Es gibt nur wenig
Situationen bei denen man darauf verzichten kann (zum
Beispiel psychiatrische Patienten). Zu den Vitalwerten zählt man

\begin{itemize}
\item die \gEs{Herzfrequenz}, 
\item den \gEs{Blutdruck} und 
\item die \gEs{Atemfrequenz}.
\end{itemize}

Die Messung der Herzfrequenz erfolgt mittels Puls fühlen oder
mittels Pulsoxymeter. Beim Puls fühlen wird für
\gEs{15~Sekunden} der Puls gefühlt und die Pulsschläge
gezählt. Anschließend wird das Ergebnis \gEs{mit vier
multipliziert}, man erhält die Pulsschläge \gEs{pro Minute}
(\nicefrac{x}{min}). Neben der Frequenz wird auch die
Regelmäßigkeit beurteilt, das heißt ob der Puls \gE{rhythmisch} oder \gE{arrhythmisch} ist.

Alternativ kann die Messung mittels des \gE{Pulsoxymeters}
erfolgen (\prettyref{topic:pulsoxymetrie}). Dabei misst ein
Gerät die Herzfrequnz und zeigt diese an. Es ist jedoch zu
beachten, dass bei Notfallpatienten die Pulsoxymetrie oft
nicht funktioniert (z.B. Zentralisation bei Schock, kalte Finger,\ldots).

Die Messung des Blutdrucks ist unter \gSee{topic:rr-messung} beschrieben.

Bei der Messung der Atemfrequenz wird die Anzahl der Atemzüge pro Minute gezählt.\index{Atemfrequenz}

% TODO Atmung






\subsection{Messung der Herzfrequenz -- HF}
\label{topic:pulsmessung}

 
\gDefinition{Als \gT{Herzfrequenz} (\gT{HF})
bezeichnet man die Anzahl der Schläge des Herzens pro Minute.} 

Sie wird mittels Fühlen des Pulses
ermittelt. \marginpar{\gE{\scriptsize Die Blutgefäße werden im Detail unter
\gSee{topic:blutgefaesse-wichtige} besprochen.}} 
Dieser wird idealerweise an der
\gEs{Radialarterie} (A. \gE{radialis}) bzw. bei nicht
ansprechbaren oder schockierten Patienten an der
\gEs{Halsschlagader} (A. \gE{carotis}) ertastet. Bei
Neugeborenen und Säuglingen kann man auch an der
\gEs{Oberarmarterie} (A. \gE{brachialis}) fühlen. Gemessen
wird immer mit zwei bis drei Fingern. Nie jedoch mit dem Daumen, da hier meistens nur der eigene Puls gefühlt wird und nicht der des
Patienten. Wenn an der Halsschlagader (A. \gE{carotis}) die
Herzfrequenz ermittelt wird, darf nie auf beiden Seiten
gleichzeitig gemessen werden, da sonst die Gefahr besteht,
dass es zu einer Unterversorgung des Gehirn durch Blut kommen
kann.

\paragraph{Durchführung}

Es wird an der Radialarterie (A. \gE{radialis}) der Puls mit
zwei bis drei Fingern ertastet. Nun werden 15 Sekunden lang
die Schläge gezählt. Multipliziert man den Messwert mit vier, erhält man  die Herzfrequenz (HF) des Patienten. Der Normalwert liegt beim Erwachsenen zwischen 60 und 100 Schlägen pro Minute.
Dokumentiert wird sie mit einem HF (z.B.: HF\,=\,70)

\subsection{Messung des Blutdrucks -- RR}
\label{topic:blutdruckmessung}

Der Blutdruck ist jener Druck, welcher in unseren Blutgefäßen
herrscht. Wenn wir vom \emph{"`Blutdruck"'} sprechen beziehen wir
uns auf den Blutdruck in den
Arterien\footnote{\emph{Arterie:} Vom Herzen wegführendes
Blutgefäß.}. Er ist in Kombination mit der Herzfrequenz
einer der wichtigsten Messwerte im
Rettungsdienst.

Im Idealfall werden zwei Werte ermittelt: Der
\gEs{systolische Wert} ("`obere Wert"') ist jener Druck im
Blutgefäßsystem, der  während der \gE{Auswurfsphase} des Herzens, in der sich
die Herzkammern anspannen, zusammenziehen und Blut auswerfen, herrscht. Der
\gEs{diastolische Wert} ("`untere Wert"') ist jener Druck im
Blutgefäßsystem, der während der \gE{Entspannungsphase} des Herzens,
in der sich die Herzkammern mit
Blut füllen, herrscht. 

Die \gEs{Einheit}, in der der Blutdruck angegeben wird, ist
\gT{Millimeter Quecksilbersäule}
(\gT{mmHg}).\index{Millimeter Quecksilbersäule}\index{mmHg}\footnote{\gE{Hg} ist
das chemische Symbol für Quecksilber.}

\gMarried{
Der Blutdruck kann auf verschiedene Arten gemessen werden. 
%Üblich ist die
%\gE{indirekte Methode}, bei welcher mittels einer Druckmanschette an einer
%Extremität gemessen wird bei welchem Druck noch ein Blutfluß stattfindet. Dies
%erlaubt Rückschlüsse auf den Druck im Blutgefäß. 

Dazu gibt es unterschiedliche Messgeräte, grundsätzlich unterscheidet man
zwischen elektronischen und rein mechanischen Messgeräten.

\gEs{Elektronische Messgeräte}  werden oft von Patienten zur Selbstkontrolle des
Blutdrucks verwendet. Diese Geräte sind jedoch vergleichsweise fehleranfällig
und sind für einen professionellen Einsatz nicht geeignet. 

Für den
professionellen Einsatz gibt es auch spezielle
elektronische Blutdruckmessgeräte die zum Beispiel in Überwachungsmonitoren zum Einsatz
kommen, doch auch diese Geräte sind noch relativ fehleranfällig wenn die
Messbedingungen nicht optimal sind. 

% Der Blutdruck kann mittels elektronischer
% Blutdruckmessgeräte ermittelt werden. Diese sind aber im
% Rettungsdienst nicht in Verwendung, da es immer wieder zu
% Gerätefehler kommen kann. Außerdem können die Batterien leer
% werden. Solche Geräte haben meistens die Patienten für die
% tägliche Überwachung ihres Blutdrucks. Leider bedienen aber
% viele die Geräte falsch und fürchten sich dann
% lebensbedrohlich krank zu sein.

In der Praxis wird die Blutdruckmessung mit \gEs{mechanischen Messgeräten}
durchgeführt. Dabei wird eine aufblasbare Manschette mit angeschlossenem Dosenmanometer und
Pumpballen mit einem Ablassventil benötigt. Im Idealfall wird auch noch ein
Stethoskop verwendet. 
}{
\begin{center}
     
  \includegraphics[width=4cm]{./images/noauth/GER-rr-elektrisch-oberarm.png}
  
   \includegraphics[width=4cm]{./images/noauth/GER-rr-elektrisch-unterarm.jpg}
   
   \includegraphics[width=4cm]{./images/noauth/GER-rr-messgeraet.jpg}
\end{center}
 }

\subsubsection{Messung nach Riva-Rocci}\marginpar{\gZ{\scriptsize Benannt nach
dem Erfinder Scipione Riva-Rocci.}}

\gDefinition{Die Blutdruckmessung nach Riva-Rocci ist eine mechanische Messung mittels einer Druckmanschette.}

Umgangssprachlich wird die Blutdruckmessung nach Riva-Rocci als
\emph{"`RR-Messung"'}, bzw. der ermittelte Wert als \emph{"`RR"'} bezeichnet.

Bei der RR-Messung gibt es zwei Möglichkeiten zu
einem Wert zu kommen. Die genauere ist die
\gEs{auskultatorische Methode} (durch hören von
Strömungsgeräuschen). Die zweite Möglichkeit ist die
\gEs{palpatorische Methode} (durch tasten des Pulses). 

\gEs{Der diastolische Wert kann nur
bei der auskultatorischen Methode ermittelt werden!} 

\paragraph{Durchführung}

\subparagraph{\gT{Auskultatorische Methode}}

Die Manschette sollte faltenfrei und nicht über Kleidung angelegt werden. Die
Mitte der Luftkammer soll über der Oberarmarterie zu liegen kommen. Manche
Manschetten haben einen Markierungspfeil für die Oberarmarterie, der Pfeil soll
dann oberhalb dieser zu liegen kommen. Man achte auf die Lage der Innen- bzw.
Außenseite!

Nun wird
der Puls an der Radialarterie (A. radialis) getastet. Die Manschette wird
solange aufgepumpt, bis der Puls nicht mehr getastet werden kann. Dann werden
noch 20-30\,mmHg weiter aufgepumpt. Das Stethoskop wird in der Ellenbeuge auf
die Oberarmarterie (A. brachialis, Verlauf eher zur Körpermitte) angelegt. Die
Luft wird langsam aus der Manschette abgelassen. 

Der Wert bei dem man
erstmalig ein Pochen \gZ{(Korotkowsche Geräusche)} hört gilt als
\gEs{systolischer Wert}. Die Luft wird weiter abgelassen. Der Wert, bei dem das
Pochen verschwindet, ist der \gEs{diastolische Wert}. \footnote{Es kann
vorkommen dass das Pochen nie verschwindet sondern nur plötzlich leiser wird.
Dann gilt der Wert bei dem das Pochen \emph{deutlich} leiser wird als
diastolischer Wert.}

Die
Werte werden auf die nächste Fünferstelle abgerundet (z.B.:
122\,/\,82\,mmHg wird abgerundet auf 120\,/\,80\,mmHg). Dokumentiert
wird der Blutdruck mit RR (z.B.: RR\,=\,125\,/\,75\,mmHg). Der
Normalwert beim Erwachsenen liegt ca. bei RR 120\,/\,80\,mmHg
(\gSee{topic:blutdruck}).


\begin{figure}[H]
    \begin{center}
   
   
    \includegraphics[width=10cm]{./images/noauth/GER-rr-geraeusche.png}

   
    \gCaptionof{figure}{\gAuthNotesPicLegalNotOk{GER-rr-geraeusche.png}{}Prinzip
    der auskultatorischen RR-Messung: In Abhängigkeit vom außen anliegenden
    Manschettendruck kommt es zu einem unterschiedlich starken Blutfluss und dadurch entstehenden, hörbaren, Strömungsgeräuschen.}
    \end{center}
\end{figure}




\subparagraph{\gT{Palpatorische Methode}}
Die Manschette sollte faltenfrei, nicht über Kleidung, in der
Mitte des Oberarms angelegt werden. Nun wird der Puls an der
Radialarterie (A. radialis) ertastet. Die Manschette
wird solange aufgepumpt bis der Puls nicht mehr getastet werden
kann. Nun noch 20\,mmHg weiter aufpumpen. Die Finger bleiben
auf der Radialarterie (A. radialis). Die Luft langsam
aus der Manschette ablassen. Sobald man wieder einen Puls fühlt,
muss man den Messwert am Dosenmanometer ablesen. Dies ist der
\gEs{systolische} Wert. 

\gE{Der diastolische Wert kann mit
dieser Methode nicht ermittelt werden.} In der
Notfallmedizin ist der systolische Wert jedoch ohnehin in der
Regel der wichtigere.

\paragraph{Was ist besonders zu beachten?}

\marginpar{\includegraphics[width=\linewidth]{./images/sign_cardive.jpg}}

Die Blutdruckmanschette soll nach Möglichkeit nicht auf der Seite
\begin{itemize}
	\item eines Shuntarmes bei Dialysepatienten
	\begin{itemize}
		\item Ein Dialyse-Shunt ist ein künstlicher Kurzschluss
		zwischen einer Armarterie und einer Armvene. Dadurch wird
		die Vene erweitert und sie pulsiert. Mann kann viel mehr
		Durchfluss für die Dialyse erreichen. Die Nahtstelle
		zwischen Vene und Arterie ist jedoch sehr sensibel, wenn
		gestaut wird kann sie aufplatzen $\rightarrow$ Gefahr
		einer massiven Shuntblutung.
	\end{itemize}
		\item eines Armödems, z.B. nach einer
		Brustamputation
		\item einer Halbseitenschwäche oder -lähmung 
\end{itemize}
angelegt werden.

\gCave{Shunt-Arme sind tabu für die Blutdruckmessung!}









\section{Apparative Untersuchungen}

\subsection{Pulsoxymetrie}
\label{topic:pulsoxymetrie}



\gMarriedNG{\gTxBeschreibung}
{
Die Pulsoxymetrie ist eine einfache apparative
Untersuchungsmethode, welche mithilfe eines
\gE{Pulsoxymeters} durchgeführt wird. 
Sie dient zur nicht-invasiven Überwachung der
\gT{Sauerstoffsättigung} (\gT{Sp\ce{O2}}) des kapillären
Blutes. Dies geschieht mittels Photorezeptoren und
Infrarotlicht, mit deren Hilfe das Verhältnis der mit
Sauerstoff beladenen roten Blutkörperchen zu den nicht
beladenen bestimmt wird. 

Als weiteren Wert zeigt das Pulsoxymeter die
\gEs{Pulsfrequenz} an.
\cite{LoefflerBiochemie7,ArbeitstechnikenAZ1}

}{
% TODO Slide fehlt.
{\gSymbolText}
}

\gLayout{2}{Durchführung}
{
Es wird ein eine Klemme mit einer besonderen Lichtquelle
und einem Sensor auf den Finger des Patienten
geklipt. Das Licht wird vom
gegenüberleigenden Sensor aufgefangen, das Gerät kann aus den daraus gewonnenen Informationen auf den Sauerstoffgehalt des Blutes und die Pulsfrequenz schließen. 

Dies setzt jedoch
voraus das der Patient in der Peripherie über eine
ausreichende Durchblutung verfügt, da das Gerät sonst keine
Informationen gewinnen kann (mangels Blut zur
Beurteilung). 
}{}{}{}{}


\gFazit{Eine ausreichende Blutversorgung der Körperstelle,
an der der Sensor angebracht ist, ist Voraussetzung für die
Pulsoxymetrie.}

Das Pulsoxymeter eignet sich sehr gut für das Monitoring,
ersetzt aber niemals das wachsame Auge auf den Patienten!

\gLayout{0}{Ermittelte Werte}
{
Wie Sauerstoffsättigung wird in Prozent
angegeben und wird mit \gT{Sp\ce{O2}} abgekürzt. Der Normalwert
liegt zwischen \gEs{96-100\,\%}. Bei Patienten mit COPD
findet man jedoch oft sehr niedrige Werte (z.B.
90\,\%), ohne dass die Patienten Beschwerden haben. }{
% TODO Slide fehlt.
% \begin{itemize}
%     \item \gZ{Point}
% \end{itemize}
\begin{itemize}
  \item Normalwert: 96-100\, \%
  \item COPD-Patient: auch Werte < 90\, \% möglich ohne
  Beschwerden
\end{itemize}
}{}{}{}


\gMarriedNG{Beurteilung}
{
Bei der Interpretation des
$SpO_2$ muss man sehr vorsichtig sein! Auch ein Normalwert bedeutet nicht, dass der Patient
genug Sauerstoff hat beziehungsweise, dass keine
Sauerstoffgabe erfolgen soll. Die Sauerstoffgabe richtet
sich immer nach dem klinischen
Zustandsbild des Patienten und dessen Diagnose, und nur in
Ausnahmefällen (z.B. Asthma bronchiale,
\prettyref{topic:asthma-bronchiale}) nach dem
Sp\ce{O2}-Wert. Der Sp\ce{O2}-Wert soll immer im Verlauf beurteilt werden. }{
% TODO Slide fehlt.
% \begin{itemize}
%     \item \gZ{Point}
% \end{itemize}
{\gSymbolText}
}

\gFazit{Behandle den Patienten und nicht das Gerät!}

Immer
zuerst auf den Patienten schauen (Symptome) und erst dann auf das Pulsoxymeter!



\gLayout{0}{\gTxGefahren{} und Fehlerquellen}
{
Eine besondere Fehlerquelle ist die
Kohlenmonoxid-(\ce{CO})-Vergiftung, bei ihr wird nicht nur
der Patient schweinchenrosa obwohl er eigentlich zu wenig
Sauerstoff hat, auch das Pulsoxymeter verwechselt
Kohlenmonoxid mit Sauerstoff und zeigt falsch hohe Werte an!

Voraussetzung für eine korrekte Messung ist eine ausreichende
Durchblutung der entsprechenden Meßstelle. Z.B. durch Kälte,
Schock oder auch eine Blutdruckmessung kann die Messung
gestört oder unmöglich gemacht werden. Weitere Hindernisse
sind z.B. kalte Finger (verringerte Durchblutung) oder
Nagellack (Licht kann den Nagellack nicht durchdringen,
bzw. dessen Farbe wird verfälscht.) Mit
Nagellack-Entferner kann der Nagel eventuell
gereinigt werden. Alternativ kann der Sensor auch
um 90\textdegree gedreht werden. 

Von der Seite einstrahlendes Licht kann ebenso das Ergebnis
verfälschen, da es sich ja um eine lichtabhängige
Meßmethode handelt.
}{
\begin{itemize}
	\item \ce{CO}-Vergiftungen (falsch hohe Wert)
	\item Minderdurchblutung der Peripherie; Entweder durch Zentralisation beim Schock oder Unterkühlung.
	\item Nagellack
	\item Künstliche Fingernägel
	\item Sonneneinstrahlung von der Seite
\end{itemize}
}{}{}{}




\subsection{Blutzuckermessung}
\label{topic:blutzucker-messung}

\gLayout{0}{Interpretation und Normalwert}
{
Der Blutzuckerspiegel wird in Milligramm pro Deziliter
(\nicefrac{mg}{dl}\gZ{, enspricht mg\,\%}) gemessen. Der
Wert ist abhängig von der letzten Mahlzeit des Patienten.
Nach der Nahrungsaufnahme steigt der Blutzuckerspiegel
stark an und fällt dann wieder auf einen "`nüchternen
Normalwert"' ab. Der normale
\gEs{Nüchtern-Blutzuckerspiegel} bewegt sich zwischen
\gEs{80\,--\,120\,\nicefrac{mg}{dl}}. Im Rettungsdienst
kann der Normalbereich jedoch großzügiger angesetzt werden,
da bei Werten zwischen 60 und 200 \nicefrac{mg}{dl} keine
Notfallsymptome auftreten.

}{
\begin{itemize}
    \item Abhängig von letzter Mahlzeit!
    \item Normalwert: 80\,--\,120\,\nicefrac{mg}{dl}
    (nüchtern)
    \item Bereich ohne Notfallsymptome:
    60\,--\,200\,\nicefrac{mg}{dl}
%     \item Notfallrelevanter Hypoglykämie: <
%     60\,\nicefrac{mg}{dl}
%     \item Notfallrelevante Hyperglykämie: >
%     200\,\nicefrac{mg}{dl}
\end{itemize}
}{}{}{}

\gMarriedNG{Material \& Personal}
{
\begin{enumerate}
  \item 1 eingeschulte Person
  \item Kontamed-Box
  \item Alkoholtupfer
  \item Trockene Tupfer
  \item Blutzuckermessgerät
  \item Blutzuckermessstreifen
  \item Lanzette
\end{enumerate}
}
{}


\gMarriedNG{Durchführung}
{
%\begin{multicols}{2}
\begin{enumerate}
  \item Finger auswählen: möglichst von der nichtdominanten
  Hand.
  \item Einstichstelle bestimmen: Am Fingerendglied,
  seitlich (Rand!), nicht in die Fingerkuppe, nicht zu nahe
  am Gelenk. Außenseite des Ringfinger bevorzugen.
  \item Stelle desinfizieren: Mit Alkoholtupfer abwischen
  und gemäß Einwirkzeit einwirken lassen.
  \item Blutzuckermessstreifen halb in das Gerät schieben,
  aber noch nicht einrasten lassen.
  \item Mit einer schnellen, beherzten Bewegung den
  Patientenfinger mit der Lanzette stechen und sofort
  herausziehen \gE{("`Rein-raus-Methode"')}.
  \item Lanzette in die \gE{Kontamed-Box} entsorgen.
  \item Ersten Bluttropfen mit trockenem Tupfer abwischen.
  \item Messstreifen ganz in das Gerät schieben, das Gerät
  schaltet sich ein und führt einen Selbsttest durch. Wenn
  im Display ein Tropfen angezeigt wird ist das Gerät
  messbereit.
  \item Den Messstreifen mit dem Gerät in den Bluttropfen
  tauchen, das Blut steigt auf\footnote{Kapillarwirkung}.
  \item Beim Piepton ist genug Blut in die Testkammer
  eingeströmt und das Gerät beginnt mit der Messung und kann
  entfernt werden. 
  \item Mit einem trockenen Tupfer auf die Einstichstelle
  drücken (lassen).
  \item Ein neuerlicher Piepton zeigt an, dass die Messung
  beendet wurde und der Wert vom Display abgelesen
  werden kann. 
  \item Messstreifen entsorgen\footnote{Entsorgung gemäß
  Abfallentsorgungsplan. \opt{ORGASBFLO}{\ASBFLO{} Der
  Meßstreifen kann in den normalen Müll entsorgt werden.}}.
\end{enumerate}
%\end{multicols}
}{

}




\subsection{Körpertemperatur}

Die Körpertemperatur kann mit unterschiedlichen Mitteln und
an unterschiedlichen Stellen bestimmt werden. 

\gMarriedNG{Messarten}
{
Mit dem
\gEs{herkömmlichen Fieberthermometer} kann die Temperatur
unter der \gE{Achsel} (\gE{axillär}) oder im \gE{Enddarm}
(\gE{rektal}, v.a. bei Kleinkindern) bestimmt werden. Die
Messung im Mund ist in der Praxis eher unüblich. Es ist bei der Messung eine
Schutzhülle zu verwenden. 

Mit dem \gEs{Ohrthermometer} kann mittels Infrarot die
Temperatur im \gE{Mittelohr} bestimmt werden. Es sind die
entsprechenden \gE{Schutzkappen} zu verwenden.

Die Temperatur im Enddarm und im Mittelohr bezeichnet man
als \gT{Körperkerntemperatur}, diese beträgt um die
37\textcelsius. Die unter der Achsel gemessene Temperatur
kann ca. \nicefrac{1}{2}\textcelsius{} niedriger sein.
}{
% TODO Slide fehlt.
% \begin{itemize}
%     \item \gZ{Point}
% \end{itemize}
{\gSymbolText}
}


\section{Körperliche Untersuchungen}

\subsection{Abtasten des Abdomens}

\gLayout{2}{Abtasten des Abdomens}
{\label{topic:untersuchung-abdomen-abtasten}
Das Abdomen kann man abtasten um die Anspannung der
Bauchdecke zu beurteilen (sie voraus), Schmerzen zu
beurteilen und zu zuordnen, Fremdkörper zu ertasten und
anderes. Man beginnt damit den Patienten zu fragen ob er ihm
auch Schmerzen hat und wo er diese hat. Die Untersuchung
beginnt man in dem Quadranten der am weitesten von den
Schmerzen entfernt ist!

Man legt nun eine Hand flach auf den Bauch des Patienten und
die andere Hand flach auf die erste. Mit den Fingern drückt
man nun vorsichtig mehrere Zentimeter auf die Bauchdecke und
beurteilt, ob der Bauch weich oder verhärtet ist. Dabei
fragt man den Patienten ob er Schmerzen empfindet. Dies
wiederholt man in allen vier Quadranten.

Weiters kann man beurteilen ob der Patient Schmerzen durch
den Druck hat (Druck Schmerz) oder ob er Schmerzen beim
loslassen bekommt
(\gT{Loslassschmerz}\index{Loslassschmerz}).
}
{
% TODO VIT: Slide fehlt!
% \begin{itemize}
%     \item \gZ{Point}
% \end{itemize}
{\gSymbolText}
}{}{}{}

\subsection{Neurocheck}

%\begin{multicols}{2}

\gLayout{1}{}{
\begin{enumerate}
  \item \gEs{Bewusstseinslage}: Für eine grobe Einschätzung
  ist folgende Differenzierung ausreichend:
  	\begin{enumerate}
        \item bewusstseinsklar
        \item bewusstseinsgetrübt
        \item bewusstlos
      \end{enumerate}
      Notärzte und Notfallsanitäter verwenden zur
      Beurteilung des Bewusstseins eine detailliertere
      Skala, die so genannte Glasgow Coma Scale.
  \item \gEs{Orientierung}:
  	\begin{enumerate}
        \item zur Person (Patient kennt seinen Namen)
        \item zum Ort (Patient weiß wo er sich befindet)
        \item zur Zeit (Patient kann das ungefähre Datum
        und die ungefähre Uhrzeit nennen)
    \end{enumerate}
  \item \gEs{Pupillen} Begutachtung und Lichtreaktion
\begin{enumerate}
  \item Pupillen gleich weit?
  \item Blickstarre (\gT{Herdblick})? \index{Herdblick}
  \item \gEs{Lichtreaktion:} In jedes Auge 2\,x leuchten und
  jeweils das gleiche und das andere Auge beobachten
  
  Beurteilung von:
  \begin{enumerate}
      \item Orientierung zu Person, Ort, Zeit und Situation.
  \item Schnelligkeit der Reaktion: prompt, verzögert?
  \item Reagieren immer beide
  Pupillen?~\gZ{\footnote{\gZ{isokor}}}
    \end{enumerate}
\end{enumerate}
  \item Lichtscheu und Sehstörungen?
  \item \gEs{Kraftprobe} grobe Prüfung an oberer und
  unterer Extremität, beurteilt wird die Kraft und die
  Seitengleichheit. 
    \begin{enumerate}
  \item Patient überkreuzt die Hände geben, bei kräftigen
  Patienten nur Zeige- und Mittelfinger.
  \item Patienten bitten kräftig zuzudrücken.
  \item Mit Handflächen von unten gegen die Zehen drücken
  und Patienten bitten \emph{"`auf das Gaspedal zu
  drücken"'}.
  \item Von oben auf die Zehen drücken und Patienten bitten
  die Zehen Richtung Kopf zu heben.
\end{enumerate}
  \item \gEs{Nackensteifigkeit} (\gT{Meningismus}) Der
  Patient liegt flach am Rücken. Der Untersucher fasst den
  Kopf des Patienten und versucht das Kinn an die Brust zu
  führen. Der Patient darf dabei nicht mithelfen.
  Normalerweise ist diese Bewegung ohne Probleme möglich.
  
  Wenn eine Nackensteifigkeit vorliegt ist die Beugung nicht
  möglich und der Versuch schmerzhaft, der Patient bewegt
  evt. den ganzen Rumpf mit. Die Bedeutung der
  Nackensteifigkeit wird unter \gSee{topic:meningitis}
  besprochen. Bei jedem Patienten mit Fieber ist auf das
  Vorliegen einer Nackensteifigkeit zu untersuchen.  
  \item \gEs{Blutzuckermessung} Eine Blutzuckermessung soll
  grundsätzlich bei einem Neurocheck erfolgen.
\end{enumerate}
}{
\begin{itemize}
  \item Bewusstseinslage
  \item Orientierung
  \item Pupillen
  \item Lichtscheue, Sehstörungen
  \item Kraftprobe
  \item Nackensteifigkeit
  \item Blutzuckermessung
\end{itemize}
}{}{}{}

%\end{multicols}

\gLayout{57}{}
{}{}
{
\begin{tabularx}{\linewidth}{XXX}
\includegraphics[width=\linewidth]{./images/motal-aass/neuro1.jpg}
&
\includegraphics[width=\linewidth]{./images/motal-aass/neuro2.jpg}
&
\includegraphics[width=\linewidth]{./images/motal-aass/neuro3b.jpg}
\\
Pupillenreaktion, Lichtscheue, Sehstörungen 
&
Kraftprobe an den Armen
& 
Kraftprobe an den Beinen
\\\addlinespace
\includegraphics[width=\linewidth]{./images/motal-aass/neuro4.jpg}
&
\includegraphics[width=\linewidth]{./images/motal-aass/neuro5.jpg}
&
\includegraphics[width=\linewidth]{./images/motal-aass/neuro6.jpg}
\\
Nackensteifigkeit
&
Nackensteifigkeit
&
Blutzuckermessung
\\\addlinespace\bottomrule
\end{tabularx}
}{Bilderserie: Neurocheck}{}

% FEHLT: Bild für Orientierung (stattdessen jetzt 2x
% Nackensteifigkeit)

\subsection{Traumacheck}\label{topic:traumacheck}
\gDefinition{Suche nach Verletzungen oder Verletzungsanzeichen am (liegenden)
Patienten.}

\gLayout{0}{Ausführlicher Traumacheck}
{
Der
hier beschriebene Traumacheck ist ein "`ausführlicher"' (vollständiger)
Traumacheck. In der Praxis werden beim Patienten zur raschen Einschätzung manche Punkte zuerst
gemacht und andere Schritte auf später verschoben (\gE{schneller} vs.
\gE{ausführlicher} Traumacheck).

\begin{itemize}
\item Patient orientiert? Unfallhergang erinnerlich? Schmerzen?

\end{itemize}
\begin{itemize}
    \item Abtasten mit festem Griff, \gEs{bei Schmerzen die vom Patienten erwähnte
    Stelle als letztes untersuchen}: von Kopf bis Fuß
    \begin{itemize}
        \item Schädel: Schädeldach, Hinterkopf,
        Stirn, Wangenknochen, Nasenbein, Kiefer.
        \item Hals.
        \item Obere Extremität: Schultern, Arme, Hände.
        \item Brustkorb: vorne und seitlich (Prellmarken? Atemmechanik?, gegen leichten
        Druck einatmen lassen).
        \item Bauch in 4 Quadranten (Prellmarken?
        Abwehrspannung/harte Bauchdecke?),
        \item Becken (stabil?), Schambein.
        \item Untere Extremität: Beine (sichtbare Verletzungen?), Füße:
        Schuhe und Socken ausziehen)
    \end{itemize}
    \item Durchblutung, Motorik, Sensibilität (\gT{DMS})
    \begin{itemize}
        \item Finger und Zehen: Kontrolle von \gE{Motorik} (bewegen lassen), \gE{Sensibilität} ({\quotedblbase}Spüren
        Sie das?{\textquotedblleft}) und \gE{Durchblutung}:
        
        \item \gT{Nagelbettprobe} (\gEs{Rekapillarisierungszeit}): leichter Druck auf
        Nagelende bis Nagelbett weiß ist, loslassen. Zeit bis Nagelbett wieder
        rosarot ist sollte nicht länger als 1-2 Sekunden sein.
        Seitenvergleich! Eine verlängerte Rekapzeit deutet auf eine
        Durchblutungsstörung hin. (einseitig oder beidseitig? lokal oder systemisch?)
    \end{itemize}
    \item Kontrolle von Mund, Nase und Ohren auf Blutungen (mit einer
    Lampe).
    
    \item \gT{Tupfertest}\index{Tupfertest}\label{topic:tupfertest}: etwas Blut aus
    Ohren oder Nase mit einem Tupfer auffangen. Sollte Liquor im Blut enthalten sein bildet sich um den roten Blutfleck
    ein heller, bernsteinfarbener Hof.
    
\end{itemize}
Bei Hinweis auf eine Schädelverletzung ist ein Neurocheck auf jeden Fall
erforderlich!
}{
\begin{itemize}
  \item Bewusstsein, Orientierung
  \item Abtasten
  \begin{itemize}
    \item Schmerzen
    \item Blutungen
    \item Fehlstellungen
    \item DMS
    \item Atemmechanik
    \item harte Bauchdecke
  \end{itemize} 
  \item Kontrolle von Mund, Nase, Ohren
  \item Tupfertest
\end{itemize}
}{}{}{}

\gLayout{57}{}
{}{}
{
\begin{tabularx}{\linewidth}{XXX}
\includegraphics[width=\linewidth]{./images/motal-aass/traumacheck1.jpg}
&
\includegraphics[width=\linewidth]{./images/motal-aass/traumacheck2.jpg}
&
\includegraphics[width=\linewidth]{./images/motal-aass/traumacheck3b.jpg}
\\
Kleidung entfernen 
&
Schuhe ausziehen (Fixierung, Stiefelgriff) 
& 
Untersuchung des Kopfs 
\\\addlinespace
\includegraphics[width=\linewidth]{./images/motal-aass/traumacheck4.jpg}
&
\includegraphics[width=\linewidth]{./images/motal-aass/traumacheck5b.jpg}
&
\includegraphics[width=\linewidth]{./images/motal-aass/traumacheck6b.jpg}
\\
In alle Körperöffnungen im Kopf leuchten, Pupillenkontrolle.
&
Neben Blut kann auch Liquor austreten (Tupfertest).
&
Untersuchung des Thorax (inkl. Atemmechanik)
\\\addlinespace
\includegraphics[width=\linewidth]{./images/motal-aass/traumacheck7b.jpg}
&
\includegraphics[width=\linewidth]{./images/motal-aass/traumacheck8b.jpg}
&
\includegraphics[width=\linewidth]{./images/motal-aass/traumacheck9b.jpg}
\\
Untersuchung des Abdomen und des Beckens
& 
Untersuchung der Extremitäten (DMS, Nagelbettprobe) 
&
Nach dem Traumacheck den entkleideten Patient vor
Wärmeverlusten schützen. 
\\\addlinespace\bottomrule
\end{tabularx}
}{Bilderserie: Traumacheck}{}

% }
% {
% {\gSymbolText}
% }

\gMarriedNG{Schneller Traumacheck}
{
Um rasch zu erkennen ob der Patient vital bedroht ist werden in der Praxis
wichtige Schritte vorgezogen.
\begin{itemize}
    \item Abtasten von

\begin{enumerate}
    \item Kopf,
    \item Hals,
    \item Thorax inkl. Atemprobe,
    \item Bauch,
    \item Becken und 
    \item Oberschenkel.
\end{enumerate} 
\item Kontrolle von Mund, Nase und Ohren auf Blutungen (mit einer
    Lampe).
\item Kann der Patient \gE{Hände} und \gE{Füße} spüren und bewegen?
% KOCH: komplettes DMS würd zu lange dauern, Schuhe ausziehen, etc. 
% GABS: Rekap-Zeit ist eh im normalen Ersteinschätzungsblock dabei.
% MOTM: POI 35: "`ja klingt verünftig."'
\end{itemize}
}
{
\begin{itemize}
    \item Rasche Ersteinschätzung.
    \item Rasches Erkennen von schweren, lebensbedrohlichen Verletzungen.
\end{itemize}
{\gSymbolText}
}




































\chapter{Angewandte Sanitätshilfe ---
\gAasRsN{RS~XIII}}
\label{chp:SAN}

\emph{Maintainer: diverse.}

% \author{Gabriel, Hochreiter, Motal}

\minitoc

\clearpage
\section{Allgemeine Tätigkeiten}

\subsection{Lagerungsarten}\label{topic:lagerungsarten}\index{Lagerungsarten}

\gLayout{57}{}
{}{}
{
\begin{tabularx}{\linewidth}{XXX}
\includegraphics[width=\linewidth]{./images/motal-aass/lagerungflachb.jpg}
&
\includegraphics[width=\linewidth]{./images/motal-aass/lagerung30gradb.jpg}
&
\includegraphics[width=\linewidth]{./images/motal-aass/lagerungoberkoerperhochb.jpg}
\\
Flachlagerung &
Lagerung mit 30\textdegree{} erhöhtem Oberkörper &
Lagerung mit stark erhöhtem Oberkörper 
\\\addlinespace
\includegraphics[width=\linewidth]{./images/motal-aass/lagerungseiteb.jpg}
&
\includegraphics[width=\linewidth]{./images/motal-aass/lagerungkollaps1b.jpg}
&
\includegraphics[width=\linewidth]{./images/motal-aass/lagerungkollaps2b.jpg}
\\
Seitenlagerung &
Lagerung mit erhöhten Beinen und Verdacht auf
Wirbelsäulenverletzung & 
Lagerung mit erhöhten Beinen
\\\addlinespace
\includegraphics[width=\linewidth]{./images/motal-aass/lagerungabdomenb.jpg}
&
\includegraphics[width=\linewidth]{./images/motal-aass/lagerungthromboseb.jpg}
&
\includegraphics[width=\linewidth]{./images/motal-aass/lagerungembolieb.jpg}
\\
Lagerung zur Entspannung der Bauchdecke &
Lagerung mit erhöhter Gliedmaße &
Lagerung mit herabhängender Gliedmaße
\\\addlinespace\bottomrule
\end{tabularx}
}{Bilderserie: Lagerungsarten}{}




\clearpage
\subsection{Hygienische Händedesinfektion}
\label{topic:hygienische-haendedesfinfektion}
% \marginpar{\includegraphics[width=4cm]{images/AASdev20090228-img102.jpg}
% 
% \includegraphics[width=4cm]{images/AASdev20090228-img101.jpg}}
\gDefinition{Desinfektion der Hände und Handgelenke mittels
einer Desinfektionslösungund einer genormten
Handwaschtechnik.}

% \gMarried{
\begin{gWideParEnv}

\begin{multicols}{2}
\paragraph{Warum?} 
Die Hände des Personals sind die
häufigsten Keimüberträger. 

Die normale Händewaschung entfernt nur groben Schmutz und
sonstige mechanische Verunreinigungen. Die
Desinfektionlösung reduziert gezielt aktiv viele Keime. 

Versuche mit Fingerfarben
zeigen, dass mit einer herkömmlichen Handwaschtechnik große
Teile der Hand ausgelassen werden (Fingerzwischenräume,
Handgelenk, Daumenkuppe,\ldots).

\begin{minipage}{\linewidth}

\paragraph{Ziel}
\begin{itemize}
    \item Keimreduktion
    \item Vermeidung von \gEs{Kreuzinfektionen}
    \begin{itemize}
        \item Patient1 ${\rightarrow}$ Sanitäter ${\rightarrow}$ Patient 2
        \item \gEs{{\textquotedblleft}Unterbrechung der
        Infektionskette{\textquotedblright}}
    \end{itemize}
\end{itemize}
\paragraph{Wann durchzuführen?}
\begin{itemize}
    \item Dienstantritt
    \item \gEs{Nach jedem Patienten!!}
\end{itemize}
\paragraph{Womit?}~
\begin{itemize}
    \item Händedesinfektionslösung:
    Sterilium\texttrademark{}, Manopronto\texttrademark{}
    o.ä.
\end{itemize}
\end{minipage}
\end{multicols}
\end{gWideParEnv}
% } {
% \begin{center}
% \includegraphics[width=3.961cm,height=3.169cm]{images/AASdev20090228-img102.jpg}
% 
% \includegraphics[width=4.001cm]{images/AASdev20090228-img101.jpg}
% \end{center}
% }

\begin{gWideParEnv}

\begin{minipage}{\linewidth}

\paragraph{Vorgehen}\vspace{1em}~ 

\begin{center}
\includegraphics[width=0.23\linewidth]{./images/hygiene/haendedesinfektion01.png}\hspace{0.01\linewidth}
\includegraphics[width=0.23\linewidth]{./images/hygiene/haendedesinfektion02.png}\hspace{0.01\linewidth}
\includegraphics[width=0.23\linewidth]{./images/hygiene/haendedesinfektion03.png}\hspace{0.01\linewidth}
\includegraphics[width=0.23\linewidth]{./images/hygiene/haendedesinfektion04.png}

\includegraphics[width=0.23\linewidth]{./images/hygiene/haendedesinfektion05.png}\hspace{0.01\linewidth}
\includegraphics[width=0.23\linewidth]{./images/hygiene/haendedesinfektion06.png}\hspace{0.01\linewidth}
\includegraphics[width=0.23\linewidth]{./images/hygiene/haendedesinfektion07.png}\hspace{0.01\linewidth}
\includegraphics[width=0.23\linewidth]{./images/hygiene/haendedesinfektion08.png}
\end{center}
\gCaptionof{figure}{Hygienische Händedesinfektion. \gSourcePicture{BEMAW.}}
\vspace{1em}
\begin{multicols}{2}

\begin{enumerate}
    \item Bei Dienstantritt: Ringe, Uhren, Armbänder,
    Schmuckstücke, etc. entfernen, Handewaschung
    durchführen, um grobe Verunreinigungen zu entfernen.
    \item Eine Portion alkoholisches
    Händedesinfektionsmittel (ca. 3\,ml = 1~Hohlhand) aus
    dem Spen\-der entnehmen.
    \item mittels Ellenbogen, vollkommene Benetzung beider Hände und
    Handgelenke
    \item Handinnenflächen gegeneinander reiben
    \item geschlossene Finger
    \item Handgelenke einreiben
    \item Mit ineinander verschränkten Fingern Handinnenflächen
    gegeneinander reiben.
    \item Mit rechter Handinnenfläche linken Handrücken reiben und
    umgekehrt, dabei Finger ineinander verschränken.
    \item Hände ineinander verhaken und Finger gegeneinander bewegen
    \item Daumen mit gegenüberliegender Hand vollständig umschließen
    und rotierend reiben.
    \item Daumenkuppen nicht vergessen!
    \item Fingerkuppen im Handteller kreisförmig einreiben bis Alkohol
    verdunstet ist
\end{enumerate}
\gEs{Einwirkzeit beachten!} (lt. Herstellervorgaben, i.d.R. mind. 30 Sekunden) 
\end{multicols}
\end{minipage}
\end{gWideParEnv}

\nocite{KirschnickPflege2}



\section{Traumatechniken}
% Michael Motal

\subsection{Helmabnahme}
\label{topic:helmabnahme}\index{Helmabnahme}

Immer, außer der Patient verweigert.

\paragraph*{Material und Personal}~

\begin{itemize}
\item 2 Helfer
\item 1 HWS-Schiene für später!
\end{itemize}

\paragraph*{Vorgehen}~

\begin{multicols}{2}
\begin{enumerate}
    \item \gEs{Pat. auffordern den Kopf nicht zu bewegen.}
    \item Ein Helfer (\gE{hinterer} Helfer) fixiert den Helm von hinten und lässt nicht mehr los,
    bis der Helm unten ist. Zwischen Helm und Schritt eine Helmhöhe
    Abstand lassen, damit der Helm in einer Bewegung abgezogen werden kann
    \item Ein Helfer (\gE{seitlicher} Helfer) spricht den Patienten von vorne an, öffnet Visier und
    Kinnriemen
    \item Der \gE{seitliche} Helfer stabilisiert den Kopf: eine Hand fasst mit
    Daumen und Zeigefinger das Unterkiefer, die andere Hand stützt den
    Kopf von unten. Während der Helm abgezogen wird rutscht diese Hand
    etwas nach
    \item Der \gE{hintere} Helfer zieht den Helm in einer S-förmigen Bewegung
    ab: den Helm zuerst kippen bis die Nasenspitze sichtbar ist, dann die
    Unterkante in der gegengleichen Bewegung zum Helfer bewegen
    \item Der \gE{hintere} Helfer zieht den Helm sanft fertig ab und legt ihn
    sicher ab
    \item Der \gE{hintere} Helfer übernimmt den Kopf im doppelten C-Griff:
    Daumen und Zeigefinger um ein Ohr, Kopf leicht auf Zug halten.
    \item Der \gE{seitliche} Helfer legt ein Stifneck an (siehe unten)
\end{enumerate}
\end{multicols}

\gLayout{57}{}
{}{}
{
\begin{tabularx}{\linewidth}{XXX}
\includegraphics[width=\linewidth]{./images/motal-aass/helm1.jpg}
&
\includegraphics[width=\linewidth]{./images/motal-aass/helm2.jpg}
&
\includegraphics[width=\linewidth]{./images/motal-aass/helm3.jpg}
\\
Während des Annäherns den Patient auffordern, den Kopf
nicht zu bewegen. 
&
Helm sofort fixieren, Visir öffnen und Patient ansprechen.
& 
Ggfs. Patient synchron und schonend auf den Rücken drehen.
\\\addlinespace
\includegraphics[width=\linewidth]{./images/motal-aass/helm4.jpg}
&
\includegraphics[width=\linewidth]{./images/motal-aass/helm5.jpg}
&
\includegraphics[width=\linewidth]{./images/motal-aass/helm6.jpg}
\\
Ggfs. Brille entfernen, Kinnriemen öffnen und behindernde
Kleidung entfernen. &
Halswirbelsäule stabilisieren.
& 
Helm kippen bis von oben die Nasenspitze sichtbar ist,
danach vorsichtig nach hinten wegziehen 
\\\addlinespace
\includegraphics[width=\linewidth]{./images/motal-aass/helm7.jpg}
&
\includegraphics[width=\linewidth]{./images/motal-aass/helm8.jpg}
&
\\
Nach dem Abziehen Stabilisierung des Kopfes übernehmen. 
& 
An die HWS-Schiene denken.
&
\\\addlinespace\bottomrule
\end{tabularx}
}{Bilderserie: Helmabnahme}{}

\subsection{Anlegen einer
HWS-Immobilisationsschiene (Stifneck\texttrademark)}
\label{topic:hws-schiene-anlegen}\index{HWS-Schiene!anlegen}\index{Stifneck|see{HWS-Schiene}}

Allgemeine
Beschreibung: \gSee{topic:hws-schiene-beschreibung}

\begin{wrapfigure}{r}{2cm}
\includegraphics[width=\linewidth]{./images/noauth/GER-stiffneck.jpg}
\end{wrapfigure}Die HWS-Schiene dient der Stabilisierung der
Halswirbelsäule. Sie soll bei jedem Verdacht auf eine
Verletzung der Halswirbelsäule eingesetzt werden; auf jeden
Fall bei jeder Helmabnahme!

\paragraph*{Material und Personal}~

\begin{itemize}
    \item 2 Helfer
    \item 1 HWS-Schiene (z.B. Stifneck\texttrademark)
\end{itemize}

\paragraph*{Vorgehen}~

\begin{multicols}{2}
\begin{enumerate}
    \item \gEs{Pat. auffordern den Kopf nicht zu bewegen.}
    \item nach Helmabnahme: Kopf weiterhin leicht auf Zug halten
    \item der hintere Helfer hält Kopf/HWS in Neutralstellung
    \item der seitliche Helfer misst die Entfernung von Schulter bis
    Unterkiefer in Fingerbreiten aus
    \item Stifneck entsprechend einstellen: gemessene Entfernung einstellen
    \item Verschiebemechanismus des Stifneck blockieren, ab nun darf kein
    Verstellen mehr möglich  sein
    \item HWS-Schiene zuerst am Kinn anlegen/anpassen und
    festhalten
    \item Klettverschlussband so einfalten, dass es nicht am Boden
    schleifen kann (hält sonst nicht mehr)
    \item Nackenteil unter Pat. durch führen (Vorsicht:
    Klettverschlussband)
    \item fest schließen
    \begin{itemize}
        \item[] Kinnspitze sollte ein kurzes Stück über Kante hinaus schauen
        (kein Durchrutschen) aber nicht allzu weit (falsch eingestellt)
    \end{itemize}
\end{enumerate}
\end{multicols}

\gLayout{57}{}
{}{}
{
\begin{tabularx}{\linewidth}{XXX}
\includegraphics[width=\linewidth]{./images/motal-aass/hwsschiene1.jpg}
&
\includegraphics[width=\linewidth]{./images/motal-aass/hwsschiene2.jpg}
&
\includegraphics[width=\linewidth]{./images/motal-aass/hwsschiene3.jpg}
\\
Nach der Helmabnahme folgt sofort das Anlegen einer
HWS-Schiene. 
&
Der Kopf bleibt fixiert, die passende Größe der HWS-Schiene
wird bestimmt und eingestellt. 
& 
HWS-Schiene zuerst am Kinn anlegen und danach den
Nackenteil unter dem Patient durch führen.
\\\addlinespace\bottomrule
\end{tabularx}
}{Bilderserie: Anlegen einer HWS-Schiene}{}

\subsection{Schaufeltrage}
\label{topic:schaufeltrage-handhabung}\index{Schaufeltrage}

Allgemeine
Beschreibung: \gSee{topic:schaufeltrage-beschreibung}

Die Schaufeltrage setzt sich aus 2 Längshälften zusammen;
die Fußteile sind längenverstellbar. Die Kopfteile haben eine fixe Länge.

Sie dient der möglichst bewegungsfreien Rettung bei Verdacht
auf Wirbelsäulen-, Becken- oder Oberschenkelverletzungen. Letztendlich soll der Patient
auf der Vakuummatratze gelagert werden. Sollte der Patient größer
als die Trage sein müssen Kopf und Oberkörper auf jeden Fall
trotzdem auf der Trage zu liegen kommen. Ausnahme: Beinverletzung, bei
der eine Wirbelsäulenverletzung ausgeschlossen werden kann.

\gZ{Alternativ kann das Spineboard angewendet werden.}

\paragraph*{Material und Personal}~

\begin{itemize}
    \item 2 Helfer
    \item Schaufeltrage
    \item Gurte für Schaufeltrage
\end{itemize}

\paragraph*{Vorgehen}~

\begin{multicols}{2}
\begin{enumerate}
    \item Schaufeltrage neben Patienten in zusammengebautem Zustand auf die
    richtige Länge einstellen. Sollte ein Patient zu groß sein
    müssen Kopf und Oberkörper auf jeden Fall trotzdem auf der Trage
    liegen! Vorsicht: Markierung der max. Länge bei Schraubmodell
    \item Öffnen der Trage: Schnappverschlüsse; je nach Methode nur den
    unteren oder beide öffnen.
    \item Aufschaufeln, 1. Methode: Fußseite offen, Pat. wird von oben
    kommend wie auf eine Schere aufgeladen, die Trage unter ihm
    geschlossen. Kontrolle des sicheren Schlusses!
    
    Diese Methode ist besonders bei Becken- und Oberschenkelverletzungen
    sinnvoll und der zweiten Methode vorzuziehen da der Patient nicht am Becken
    bewegt wird.
    \item Aufschaufeln: 2. Methode: beide Teile getrennt, Pat. leicht an
    Schulter und Hüfte zur Seite gedreht und jew. Tragenteil
    daruntergeschoben. Hälften zuerst am Kopf-, dann am Fußende
    schließen. Kontrolle des sicheren Schlusses!
    \item Angurten: 2 Gurte; jeden Gurt zuerst um den Handgriff der einen,
    dann der anderen Hälfte führen, dann über dem Patienten
    schließen und festziehen.
    \item Umlagern auf Vakuummatratze: Trage mit vorbereiteter Matratze
    (\gSee{topic:vakuummatratze-anwendung}) neben Patienten stellen um diesen
    einen möglichst geringen Weg tragen/bewegen zu müssen.
    \item Gurte öffnen und entfernen.
    \item Hälften trennen: beide Schnappverschlüsse öffnen
    \item Hälften parallel zum Patienten flach wegziehen. Der Patient soll
    hierbei nicht bewegt werden!
    \item Sollten die Schnappschlösser klemmen kann man die Trage wie bei
    der Scherentechnik an nur einer Seite öffnen und in V-Form (gleichzeitig, je
    ein Helfer auf einer Seite) entfernen.
\end{enumerate}
\end{multicols}



\subsection{Vakuummatratze}
\label{topic:vakuummatratze-handhabung}\index{Vakuummatratze!Handhabung}

Allgemeine
Beschreibung: \gSee{topic:vakuummatratze-beschreibung}

zur Ruhigstellung des gesamten Körpers, bei Verdacht auf mögliche
Wirbelsäulenverletzungen, Becken- oder Oberschenkelfrakturen %(lockere
%Indikationsstellung!)

Am Ende soll die Matratze an die Körperform der Patienten eng
angepasst sein und somit ein Bewegen verhindern.

\paragraph*{Material und Personal}~

\begin{itemize}
    \item 2 Helfer
    \item Fahrtrage
    \item Vakuummatratze
    \item Accuvac
    \item Leintuch
    \item ev. Adapter, je nach Matratzentyp
\end{itemize}


\paragraph*{Vorgehen}~

\begin{multicols}{2}
\begin{enumerate}
    \item Fahrtrage auf tiefste Stufe stellen (ganz unten)
    \item Vorbereiten: Vakuummatratze auf Fahrtrage legen. Die weichere
    Seite ist die Patientenseite, Ventil beim Kopf. Kugeln so verteilen,
    dass der gesamte Körper gut, sowie die betroffene Stelle besonders gut
    fixiert werden kann.
    \item Vorabsaugen mit Accuvac bis die Matratze gut formbar aber nicht allzu
    fest ist.
    \item Vorformen, Ventil verschließen. Wenn vorhanden Leintuch auf
    die Matratze legen.
    \item Angegurteten Patienten mit Schaufeltrage auf die Vakuummatratze
    legen, Schaufeltrage entfernen (\gSee{topic:schaufeltrage-handhabung})
    \item Absaugen mit Accuvac, dabei zügig Matratze an den Patienten
    anpassen. Besonderes Augenmerk gilt der Fixierung der verletzten
    Regionen sowie dem Kopf ({\quotedblbase}Ohren{\textquotedblleft}
    formen: Ecken neben dem Kopf hoch formen).
    \item Ist kein weiteres Absaugen mehr möglich und die Matratze gut
    angepasst (Matratze muss jetzt bretthart sein) Ventil schließen,
    Accuvac abdrehen.
    \item Pat. an Fahrtrage angurten.
\end{enumerate}
\end{multicols}

\subsection{Spineboard}
\label{topic:spineboard-anwendung}\index{Spineboard!Anwendung|gZ}

Allgemeine Beschreibung:
\gSee{topic:spineboard-beschreibung}

\ASBFLO{Das Spineboard gehört derzeit nicht zur
Standardausstattung und wird nicht routinemäßig verwendet.}

% \gTakeNotesMedium

\clearpage
\subsection{Rettungskorsett (KED-System)}
\label{topic:rettungskorsett-anwendung}\index{Rettungskorsett!Anwendung}\index{KED-System!anlegen}

Allgemeine
Beschreibung: \gSee{topic:rettungskorsett-beschreibung}

\begin{wrapfigure}{r}{3cm}
\includegraphics[width=\linewidth]{./images/geraete/KED.jpg}
\end{wrapfigure}Zur schnellen Rettung sitzender oder schwer zugänglicher
Patienten; wirbelsäulenstabilisierend. Am Kopfgriff des
Korsetts kann, falls nötig, ein Kranhaken eingehängt
werden. Das KED wird immer gemeinsam mit einem Stifneck
verwendet.

\paragraph*{Material und Personal}~

\begin{itemize}
    \item 2-3 Helfer
    \item Rettungskorsett
    \item HWS-Schiene
    \item vorbereitete Vakuummatratze
\end{itemize}


\paragraph*{Vorgehen}~

\begin{multicols}{2}
\begin{enumerate}
    \item Patient aufsuchen und sofort HWS-Schiene anlegen
    (\gSee{topic:hws-schiene-anlegen})
    \item Ein Helfer stabilisiert von nun an am besten von hinten den Kopf um seitliche Bewegungen des Kopfes
    zu vermeiden.
    \item KED \gE{vorsichtig} in aufgefaltenem Zustand
    hinter dem Patienten platzieren. Die Oberkanten der
    Flügel sollen in der Achsel des Patienten zu liegen
    kommen.
    \item Die Beingurte vom Korsett lösen und baumeln
    lassen.
    \item der mittlere und untere Bauchgurt wird
    geschlossen und halbfest angezogen. Das Korsett soll in der Höhe nicht mehr
    verrutschen können.
    \item Kopf mit der zum System gehörigen faltbaren
    Schaumstoffunterlage so unterpolstern, dass die HWS in
    Neutralposition steht.
    \item der Kopf wird mit den Klettbändern fixiert: ein
    Band an der Stirn, das andere unter dem Kinn die HWS-Schiene
    entlang. Die Klettbänder werden am KED
    fixiert. Der Kopf soll nun seitlich nicht mehr geneigt
    werden können. Der Patient wird \gE{unter keinen
    Umständen zu einem Versuch aufgefordert!}
    \item die Beingurte werden geschlossen: von unten um
    den Oberschenkel herum nach oben; vorsichtig hin- und
    zurückziehend so eng es geht anlegen und festzurren.
    Vorsicht auf den Genitalbereich!
    \item die halbfest geschlossenen Bauchgurte nun
    festzurren
    \item den letzten, obersten Bauchgurt schließen und
    festzurren. Dieser wird erst jetzt geschlossen um die
    Atemarbeit des Patienten nicht zu behindern.
    \item der Patient ist nun entlang der gesamten
    Wirbelsäule geschient und kann mittels der 3 am Korsett
    angebrachten Griffe gehoben werden.
    \item sanftes Umlagern auf die vorbereitete
    Vakuummatratze (\gSee{topic:vakuummatratze-handhabung})
    \item Beingurte leicht lockern
    \item den obersten Bauchgurt öffnen.
    \item Vakuummatratze anformen und sichern 
    (\gSee{topic:vakuummatratze-handhabung})
       
\end{enumerate}
\end{multicols}



\subsection{Vakuum-Beinschiene}
\label{topic:vacuumbeinschiene-anwendung}\index{Vakuumbeinschiene!Anwendung}

Allgemeine
Beschreibung: \gSee{topic:vacuumbeinschiene-beschreibung}

\begin{wrapfigure}{r}{2cm}
\includegraphics[width=\linewidth]{./images/geraete/vakuumbeinschiene.jpg}
\end{wrapfigure}Schienung einer Fuß-, Knie- oder Unterschenkelverletzung. Achtung:
Oberschenkel: Vakuummatratze

1 Helfer seitlich, 1 Helfer beim Fußende


\paragraph*{Material und Personal}~

\begin{itemize}
    \item 2 Helfer
    \item Vakuumbeinschiene
    \item Accuvac
    \item ev. Adapter
\end{itemize}


\paragraph*{Vorgehen}~

\begin{multicols}{2}
\begin{enumerate}
    \item Länge anpassen, gegebenenfalls oben umfalten.
    Klettverschlussbänder herrichten.
    \item Kugeln gut verteilen. Auch am Fußende müssen genug Kugeln
    sein, um dem Fuß sowohl unter der Sohle als auch am Fußrücken
    gut Halt geben zu können
    \item Schiene leicht vorabsaugen, damit Kugeln nicht mehr verrutschen
    können
    \item Schuh am gesunden Bein ausziehen: Finden der schonensten Methode.
    Ein Helfer hält das Bein auf Zug, der andere öffnet den Schuh
    (Bänder ausfädeln, zur Not aufschneiden)
    \item vorsichtiges Wiederholen zu zweit am verletztem Bein, Socken
    ausziehen (DMS)
    \item Stiefelgriff: gleich nach Ausziehen von Schuh/Socken nimmt der
    Helfer am Fußende durch das Loch in der Schiene die Ferse des
    verletzten Fußes, die andere Hand hält den Fußrücken (Stiefelgriff).
    \gEs{Das Bein wird so auf dem Boden liegend auf Zug gehalten, der Helfer
    verbleibt so, bis die Schiene fest sitzt.}
    \item Der seitliche Helfer richtet die Schiene vorsichtig unter dem
    Patienten aus und fixiert die Schiene mit den Klettverschlussbändern.
    Um den Fuß werden Ohren gebogen und kreuzweise festgeklettet, alternativ ein
    Band über den Fußrücken und eines unter der Sohle. Der Fuß soll von allen
    Seiten so gut wie möglich umschlossen sein!
    \item Der seitliche Helfer saugt ab und zieht dabei die Klettbänder
    immer wieder fest nach. \gEs{Der andere Helfer hält nach wie vor den
    Fuß im Stiefelgriff.}
    \item Fertig abgesaugt: Ventil schließen, Accuvac abschalten. Der
    Stiefelgriff kann jetzt losgelassen werden.
\end{enumerate}
\end{multicols}

Die Schiene soll bretthart sein, das Bein leicht auf Zug fixiert. Der
Fuß soll weder seitlich noch zum Körper hin/vom Körper weg beweglich sein.



\subsection{Sam-Splint}
\label{topic:sam-splint-anwendung}\index{Sam-Splint!Anwendung}

Allgemeine
Beschreibung: \gSee{topic:sam-splint-beschreibung}

Unterarmschienung

\paragraph*{Material und Personal}~

\begin{itemize}
    \item 2 Helfer
    \item Sam-Splint
    \item Peha-Haft, Mullbinde + Leukoplast oder 2-3 Dreiecktuchkrawatten
    \item ein weiteres Dreiecktuch als Armtragetuch
\end{itemize}


\paragraph*{Vorgehen}~

\begin{enumerate}
    \item Sam-Splint in der Hälfte falten, an gesundem Arm der Länge
    nach anpassen, Knick beim Ellbogen
    \item Anlegen an verletztem Arm, sanft anpassen. Fixierung mit
    Peha-Haft, Mullbinden oder Dreiecktuchkrawatten
    \item Fingerspitzen müssen sichtbar sein (DMS), gegebenenfalls Enden
    umfalten.
    \item Lagerung des geschienten Armes mittels Armtragetuch
\end{enumerate}



\subsection{Druckverband}
\label{topic:druckverband-anlegen}\index{Druckverband!anlegen}

Der Druckverband dient der Stillung einer starken, d.h.
potentiell lebensbedrohlichen, Blutung.

\paragraph*{Material und Personal}~

\begin{itemize}
    \item 1-2 Helfer
    \item sterile Wundauflage je nach Wundgröße
    \item Druckkörper: zB Mullbinde
    \item Dreiecktuchkrawatte
\end{itemize}

\paragraph*{Vorgehen}~

\begin{multicols}{2}
\begin{enumerate}
    \item Sofort bei Erkennen einer \gE{starken} Blutung durch
    direkten Druck auf die Wunde und Hochhalten die Blutung verlangsamen (Pat. kann dies je
    nach Zustand selbst tun!)
    \item sterile Wundauflage auf die Wunde legen
    \item Druckkörper dem Wundverlauf folgend über die sterile
    Wundauflage legen, weiterhin drücken.
    \item Druckkörper mit einer Dreiecktuchkrawatte fest und möglichst
    vollständig umschließen
    \item Dreiecktuchkrawatte festziehen und mit einem festen Knoten bei
    ausreichendem Druck direkt über der Wunde fixieren
    \item 2. Knoten zur Sicherung des Ersten
\end{enumerate}
\end{multicols}

Sollte die Blutung zu einem späteren Zeitpunkt wieder
stärker werden oder der Druckverband zu locker angelegt
sein wird ein \gEs{weiterer} Druckkörper mit einer
weiteren Dreiecktuchkrawatte noch fester angelegt. \gEs{Ein
schon angelegten Druckverband wird nicht wieder geöffnet!}



% \subsection{Bilderserie}
% 
% \clearpage
% 
% \begin{gWideParEnv}
% \begin{tabularx}{\linewidth}{XXXX}
% \toprule
% \includegraphics[width=\linewidth]{./images/khd/gef-wind-rauch.jpg}
% 
% Text 1
% &
% \includegraphics[width=\linewidth]{./images/khd/gef-wind-rauch.jpg}
% 
% Text 2
% &
% \includegraphics[width=\linewidth]{./images/khd/gef-wind-rauch.jpg}
% 
% Text 3
% &
% \includegraphics[width=\linewidth]{./images/khd/gef-wind-rauch.jpg}
% 
% Text 4
% \\\midrule
% \includegraphics[width=\linewidth]{./images/khd/gef-wind-rauch.jpg}
% 
% Text 1
% &
% \includegraphics[width=\linewidth]{./images/khd/gef-wind-rauch.jpg}
% 
% Text 2
% &
% \includegraphics[width=\linewidth]{./images/khd/gef-wind-rauch.jpg}
% 
% Text 3
% &
% \includegraphics[width=\linewidth]{./images/khd/gef-wind-rauch.jpg}
% 
% Text 4
% \\\midrule
% \includegraphics[width=\linewidth]{./images/khd/gef-wind-rauch.jpg}
% 
% Text 1
% &
% \includegraphics[width=\linewidth]{./images/khd/gef-wind-rauch.jpg}
% 
% Text 2
% &
% \includegraphics[width=\linewidth]{./images/khd/gef-wind-rauch.jpg}
% 
% Text 3
% &
% \includegraphics[width=\linewidth]{./images/khd/gef-wind-rauch.jpg}
% 
% Text 4
% \\\midrule
% \includegraphics[width=\linewidth]{./images/khd/gef-wind-rauch.jpg}
% 
% Text 1
% &
% \includegraphics[width=\linewidth]{./images/khd/gef-wind-rauch.jpg}
% 
% Text 2
% &
% \includegraphics[width=\linewidth]{./images/khd/gef-wind-rauch.jpg}
% 
% Text 3
% &
% \includegraphics[width=\linewidth]{./images/khd/gef-wind-rauch.jpg}
% 
% Text 4
% \\\midrule
% \includegraphics[width=\linewidth]{./images/khd/gef-wind-rauch.jpg}
% 
% Text 1
% &
% \includegraphics[width=\linewidth]{./images/khd/gef-wind-rauch.jpg}
% 
% Text 2
% &
% \includegraphics[width=\linewidth]{./images/khd/gef-wind-rauch.jpg}
% 
% Text 3
% &
% \includegraphics[width=\linewidth]{./images/khd/gef-wind-rauch.jpg}
% 
% Text 4
% \\\midrule
% \includegraphics[width=\linewidth]{./images/khd/gef-wind-rauch.jpg}
% 
% Text 1
% &
% \includegraphics[width=\linewidth]{./images/khd/gef-wind-rauch.jpg}
% 
% Text 2
% &
% \includegraphics[width=\linewidth]{./images/khd/gef-wind-rauch.jpg}
% 
% Text 3
% &
% \includegraphics[width=\linewidth]{./images/khd/gef-wind-rauch.jpg}
% 
% Text 4
% \\\midrule
% \\\bottomrule
% \end{tabularx}
% \end{gWideParEnv}

\section{Assistenztätigkeiten}

\begin{tabularx}{\linewidth}[H]{XX}
Teile des Textes und der Handlungsanleitungen entstammen dem Skriptum
\emph{{\quotedblbase}Ärztliche
Grundfertigkeiten{\textquotedblleft}} der \emph{Besonderen
Einrichtung zur medizinischen Aus- und Weiterbildung --
Medizinische Universität Wien, Abt. Methodik und Entwicklung}, welches in Zusammenarbeit  mit dem Klinischem Institut für Hygiene und
Medizinische Mikrobiologie entstanden ist. 

Wir danken Hrn. ao.Univ.Prof. Dr. Michael Schmidts und
seinem Team für die Unterstützung!
 &
{\sffamily\itshape Powered by}\vspace{1em}

\includegraphics[width=4.551cm]{icons/muw.png}
\\
\end{tabularx}





\subsection{Intubation: Vorbereitung und Assistenz}
\label{topic:intubation-assistenz}

\paragraph*{Material}~

\begin{multicols}{2}
\begin{compactenum}
    \item Beatmungsbeutel inkl. passender Beatmungsmaske, Bakterienfilter und
    Sauerstoffreservoir
    \item \ce{O2}-Berieselungseinheit und \ce{O2}-Line
    \item Evt. Beatmungsgerät
    \item Absaugeeinheit: \par\gCa{Keine Intubation ohne
    Absaugung!}
    \item Laryngoskop, bestehnd aus
    \begin{compactenum}
        \item Spatel: gebogen oder gerade; Größen 0--4
        \item Griff
    \end{compactenum}
    \item (Endotracheal-)Tubus; verschiedene Größen
    \item Führungsdraht (\gT{Mandrin})
    \item Silikonspray
    \item Evt. Magillzange
    \item Mit Luft gefüllte Spritze zum Cuffen (10\,ml)
    \item Beißschutz: Beißkeil oder Guedel-Tubus
    \item Fixationsmaterial (Mullbinde)
    \item Stethoskop
\end{compactenum}
\includegraphics[width=\linewidth]{./images/intubation_zubehoer_op.jpg}
\gCaptionofDiff{figure}{Intubationszubehör, angerichtet auf
einem Anästhesiegerät. Was fehlt?\gMarginBugfix{\gTakeNotesLarge}
\gSourcePicture{Gabriel}
\gAuthNotesPicLegalOk{intubation\,zubehoer\,op.jpg}{cc-by}}{Intubationszubehör.}
\end{multicols}


\paragraph*{Vorbereitung}~

\begin{multicols}{2}


\begin{enumerate}
    \item Gewünschte Spatelgröße erfragen
    \item Gewünschte Tubusgröße erfragen
    \item Laryngoskop (Spatel und Griff) zusammenbauen
    \item Tubusverpackung aufschälen, Tubus aber in der Verpackung belassen
    \item Führungsdraht (Mandrin) einführen. Die Spitze muss bis zum Tubusende
    vorgeschoben vorgeschoben werden, darf aber nicht herausragen.
    \item Cuff mit Silikonspray einsprühen
    \item Stethoskop: Bei Doppelkopfstethoskopen Membran-Seite einstellen.
    \item Absaugung bereit machen:
    \begin{enumerate}
        \item Passenden Absaugkatheter wählen
        \item Verpackung teilweise aufschälen um Anschluss freizulegen, Katheter
        in der Verpackung belassen.
        \item Katheter anschließen.
        \item Funktionskontrolle: Gerät einschalten. Saugstärke?
        Batteriewarnung?
\item Absaugeinheit muss in Griffweite der Assistenten positioniert werden! 
    \end{enumerate}
    \item Beatmungsbeutel zusammenbauen und an \ce{O2}-Beriselungseinheit
    anschließen
    \item Restliches Material in Nierentasse bereitlegen
\end{enumerate}
\end{multicols}

\paragraph*{Assistenz}~
\begin{multicols}{2}


\begin{enumerate}
    \item Überprüfen ob alle Materialen bereit sind.
    \item Sauerstoffberiselung einschalten und auf 15\gLMin einstellen,
    Reservoir des Beatmungsbeutels füllen. 
    \item Durchführenden melden, dass alle Materialen bereit sind.
    \item \gZ{Evt. werden jetzt von einem Arzt Medikamente verabreicht die
    die Intubation ermöglichen. Der Durchführende wird nun einige Zeit den
    Patienten mittels Beutel-/Masken-Beatmung beatmen (sofern nicht bereits
    geschehen).}
    \item Laryngoskop in die \gEs{linke Hand} zureichen\footnote{Es ist egal ob
    der Durchführende Rechts- oder Linkshänder ist, das Laryngoskop wird
    \gE{immer in der linken Hand} gehalten.}
    \item Auf Anweisung \emph{"`Kehlkopfdruck"', "`Krikoiddruck"'} auf den
    Kehlkopf drücken.
    \item Tubus in die Rechte Hand zureichen.
    \item Auf Anweisung saugen.
    \item Beatmungsmaske vom Beatmunsbeutel trennen.
    \item \gZ{Durchführender intubiert und entfernt Führungsdraht (Mandrin)}
    \item Mit luftgefüllter Spritze Cuff aufblasen (\gE{"`cuffen"'})
    \item Beatmungsbeutel an Tubus anschließen
    \item Notarzt Stethoskop in die Ohren geben. 
    \item Es werden 3 Atemhübe verabreicht, wobei das
    Stethoskop über über dem Magen, linkem Lungenflügel und
    rechtem Lungenflügel positioniert wird. Es wird dabei überprüft ob der Tubus richtig liegt. 
    \item Der Durchführende übernimmt die Beatmung. 
    \item Einführen des Beißkeils oder des Guedel-Tubus
    \item Fixierung des Beißschutzes und des Tubus mittels Mullbinde
    \item \gZ{Ggfs. Beatmungsgerät einstellen und
    anschließen}
\end{enumerate}
\end{multicols}

\gTakeNotesMedium



\clearpage
\subsection{Aufziehen eines Medikamentes aus einer Ampulle}

Ampulle mit rotem Punkt: Sollbruchstelle, Daumen muss beim aufbrechen
auf dem roten Punkt sein. 

Ampulle mit {\quotedblbase}Halsring{\textquotedblleft}: Ampulle kann
nach allen Richtungen aufgebrochen werden. 

\marginpar{\gH{Warum mach{\textquotesingle} ich mir wegen
der Verletzungsgefahr Sorgen? Ist ja nur ein kleiner
Schnitt, ist ja meine Sache {\dots}}

Nein! Unter Normalbedingungen, wenn ich in aller Ruhe mein Medi
aufziehen kann ist es tatsächlich mein eigenes Bier. Was aber wenn
ich mich bei einer Reanimation an der Ampulle schneide? Die anderen
Kollegen sind beschäftigt (Drücken, blasen, Defi {\dots}) und dann
steh ich da mit meinem kleinen Schnitt und blute alles an. Und der
Notarzt braucht dringend das Adrenalin {\dots} Was jetzt?

Inakzeptabel!}

\paragraph{Material}~ 

\begin{itemize}
    \item Spritze
    \item Kanüle(n)
    \item Ampulle mit Medikament
    \item Trockene Tupfer
\end{itemize}

\paragraph{Aufziehen des Medikamentes}

\begin{multicols}{2}

\begin{enumerate}
	\item Arzt verlangt ein Medikament
	\item Sanitäter sucht das Medikament und zeigt die noch verschlossene Ampulle dem Arzt bevor sie geöffnet wird
    \item \gZ{Kontrolle Ablaufdatum}
    \item Kontrolle Inhalt: Verfärbung, Ausflockung?
    \item Öffnen der Ampulle
    \begin{itemize}
        \item  Glasampulle
        \begin{itemize}
            \item Ampullenspitze beklopfen, darauf achten, dass sich keine
            Flüssigkeit im Apullenkopf mehr befindet
            \item Öffnen der Medikamenten-Ampulle mit Tupfer, roter Punkt weist
            zum Körper
\marginpar{\includegraphics[width=\linewidth]{./images/ampulle_beklopfen.png}

\gCaptionof{figure}{Ampullenkopf
beklopfen. \emph{Bild: BEMAW}}}
        \end{itemize}
        \item Plastikampulle,
        {\quotedblbase}Dreh-und-Drink{\textquotedblleft}-Verschluss
        \begin{itemize}
            \item Den Verschlussknopf ohne Kontamination des Ampullenhalses
            (Sollbruchstelle) abdrehen
        \end{itemize}
    \end{itemize}
    \item Ampulle abstellen
    \item Verpackung der Spritze aufschälen, Spritze in Verpackung lassen
    \item Verpackung der Nadel aufschälen, Spritze auf Konus aufstecken
    \item Entfernen der Kanüleabdeckung (nicht abdrehen, sondern
    abziehen!)
    \item Einführen der Nadel in die Ampulle
    \begin{itemize}
        \item Beim Einführen der Kanüle in die Ampulle den Ampullenrand
        nicht berühren
    \end{itemize}
    \item Aufziehen des Medikamentes
    \item Vorsichtig entlüften
    \item Variante {\quotedblbase}Verabreichung über
    Venflon{\textquotedblleft}
    \begin{itemize}
        \item Nadel händisch abziehen und in den Kontamed verwerfen.
        \item Arzt das Medikament übergeben, Ampulle zeigen.
        \item Wird das Medikamen nicht sofort verwendet spritze mit roten
        Stoppel verschließen und Ampulle ankleben.
    \end{itemize}
    \item Variante {\quotedblbase}Zuspritzen in eine
    Infusionsflasche{\textquotedblleft}
    \begin{itemize}
        \item Gummimembran desinfizieren
        \item Kanüle nicht bis zum Schaft einstechen
        \item Medikament zuspritzen
        \item Beschriften der Flasche
    \end{itemize}
    \item Variante {\quotedblbase}direktes Spritzen{\textquotedblleft}
    \begin{itemize}
        \item Arzt nach Nadelgröße fragen
        \item Aufziehnadel händisch entfernen
        \item Spritznadel aufstecken
        \item Arzt Spritze mit Verschlusskappe reichen, Ampulle zeigen.
    \end{itemize}
\end{enumerate}
\end{multicols}

\subsection{Vorbereitung einer Infusion}

\paragraph*{Material}~

\gMarried{
\begin{itemize}
    \item Infusionsflasche
    \item Infusionsbesteck
    \item Aufhängevorrichtung
    \item alkoholhaltiges Desinfektionsmittel
    \item einige saubere Tupfer
    \item Permanentmarker
\end{itemize}
}{
\centering\includegraphics[width=0.5\linewidth]{./images/infusion_tropfkammer.jpg}
}


\paragraph*{Handlungsschritte}~

\begin{multicols}{2}
\begin{enumerate}
    \item Aufhängevorrichtung auf Flasche aufsetzen wenn nötig
    \item Eröffnen der Flaschenabdeckung (Membranschutz)
    \item Wischdesinfektion der Gummimembran
    \begin{itemize}
        \item Alkoholgetränkter Tupfer kann während weiterer Vorbereitung
        auf Membran belassen werden
    \end{itemize}
    \item Montage des Infusionsbestecks
    \begin{itemize}
        \item \gEs{Radklemme schließen}
        \item Einstich des Tropfkammerdorns in Infusionsflasche mit
        geschlossenem Lüftungsventil und geschlossener Radklemme
        \item {\quotedblbase}No touch technique{\textquotedblleft} -- Dorn der
        Tropfkammer und Gummimembran der Infusionsflasche dürfen nicht
        berührt werden
    \end{itemize}
    \item Infusionsschlauch füllen
    \begin{itemize}
        \item Flasche wenden und hochhalten
        \item Tropfkammerluftventil öffnen
        \item Tropfkammer durch Zusammendrücken bis zur Hälfte füllen
        \item Konusabdeckung des Infusionsschlauchs am anderen Ende muss bei der
        Belüftung nicht entfernt werden
        \item Infusionsflasche und Ende des Infusionsschlauchs mit nicht
        dominanter Hand halten  
        \item Radklemme öffnen
        \item Infusionsschlauch entlüften
        \item Radklemme schließen
    \end{itemize}
\end{enumerate}
\end{multicols}


\subsection{Venenverweilkanüle}
({\quotedblbase}Venflon{\textquotedblleft}): Vorbereitung und Assistenz 

\begin{figure}[H]
\includegraphics[width=\linewidth]{./images/venflon.png}

\gCaptionof{figure}{Eine periphere Venenverweilkanüle
(Venflon\texttrademark), in ihre Einzelbestandteile zerlegt.
\gSourcePicture{BEMAW}}
\end{figure}

\begin{table}[H]
\gCaptionof{figure}{Verschiedene Größen von peripheren
Venenverweilkanülen. \gSource{BEMAW (modifiziert)}}

\begin{tabular}{m{1.5619999cm}cccll}
\multicolumn{2}{l}{Nur zur Illustration!} 
&
Länge
&
 Durchfluss 
 &
 Farbcode
 &
\\
 {\O} mm & G & mm  & [\nicefrac{ml}{min}] & &
\\\midrule
\gZ{0,5 x 0,8} & \gZ{22} & \gZ{25} & \gZ{25} & blau & Sehr
dünn
\\
\gZ{0,7 x 1,0} & \gZ{20} & \gZ{32} & \gZ{55} & rosa & dünn
\\
\gZ{0,9 x 1,2} & \gZ{18} & \gZ{40} & \gZ{90} & grün &
normal
\\
\gZ{1,1 x 1,4} & \gZ{17} & \gZ{42} & \gZ{135} & weiß &
dick
\\
\gZ{1,3 x 1,7} & \gZ{16} & \gZ{42} & \gZ{170} & grau & Sehr
dick
\\
\gZ{1,6 x 2,1} & \gZ{14} & \gZ{42} & \gZ{265} & orange &
Pferd \\\bottomrule
\end{tabular}
\end{table}


\paragraph*{Material}~

\begin{multicols}{2}
\begin{itemize}
    \item Nierentasse
    \item Staubinde
    \item Venflon:
    \begin{itemize}
        \item Die Größe (= Farbe) ist vom Durchführenden vorher zu erfragen!
        \item Die Flügel werden hinunter geklappt
        \item Die farbige Lasche des Zuspritzventils in rechten Winkel zur
        Venflonachse stellen.
    \end{itemize}
    \item Einmalverschlusskappe ({\quotedblbase}roter
    Stoppel{\textquotedblleft})
    \item Fixierungspflaster
    \item einige saubere trockene Tupfer
    \item Alko-Tupfer
    \item (Einwegspritze mit 5\,ml Spülflüssigkeit)
    \item (Faserschreiber)
    \item Kontamed-Box
\end{itemize}
\end{multicols}

\paragraph*{Assistenz}~

\begin{multicols}{2}

\begin{itemize}
    \item Stauschlauch zureichen
    \item geöffnete Verpackung des Alko-Tupfer hinreichen, sodass dieser
    entnommen werden kann
    \item Kontamed in Reichweite des  Durchführenden
    bereitstellen, sodass er später die Nadel abwerfen kann.
    \item Venflon zureichen
\end{itemize}
\begin{itemize}
    \item Venflon an der Schutzkappe halten, sodass der
    Durchführende den Venflon abziehen kann.
    
\end{itemize}
\begin{itemize}
    \item Trockene Tupfer zureichen.
    \item Je nach Situation Infusionsschlauch, roten Stoppel oder gar nichts
    zureichen.
    \item Fixierpflaster zureichen
\end{itemize}

\end{multicols}

\gA{Butterfly gehört erklärt}
% TODO Butterfly



